% Folder 149, batch 02 — archivists' pages 21–40.
% Pass: Opus 5.5 (claude-opus-5-5), 2026-10-03 — first pass, unchecked against the pages by a human.
\documentclass[11pt,a4paper]{article}
\input{../preamble/grothendieck}

\folder{149}
\batch{2}
\pages{21}{40}
\foldertitle{[Autour de La "Longue Marche" à travers la théorie de Galois, pages 1 à 67] : notes manuscrites (s.d.).}
\dating{s.d. — le groupe « Autour de La "Longue Marche" à travers la théorie de Galois » (141 à 149) est daté [à partir de 1978]-1983}
\watermark{Édition de démonstration}

\begin{document}

\page{21}
\note{le théorème s'appuie sur des notations ($\Pi_K$, $\Pi_{\mathbb{Q}}$, $\mathcal{E}_{\overline{K}/K}$) introduites avant ce lot ; la page précédant la page 11 de l'auteur n'y porte pas de numéro visible.}

\textbf{Théorème 2.} Le foncteur $K \mapsto \Pi_K/\Pi_{\mathbb{Q}}$ \uncertain{des catégories des t.f.} de $\mathbb{Q}$ vers les groupoïdes profinis sur $\Pi_{\mathbb{Q}}$ est \emph{fidèle}.

\marginal{En fait, ce théorème n'est pas spécial à $\mathbb{Q}$ : il \uncertain{s'énoncerait} sur un corps de base quelconque \uncertain{$k$}, et \ill{} \uncertain{fait la diff.} \ill{} \uncertain{bien que l'énoncé} \ill{} \uncertain{fait sur} \ill{} \uncertain{l'énoncé est} \ill{} \uncertain{arith.} \ill{} \uncertain{géométrique} \ill{} (\uncertain{bien que} \ill{}) \ill{} \uncertain{mais} \ill{}.}
\note{note marginale écrite en oblique dans la marge gauche, en partie à une autre encre ; seul le début se lit sûrement.}

I.e. si deux hom. $K \overset{f,g}{\rightrightarrows} K'$ définissent des hom. de groupoïdes \add{sur $\Pi_{\mathbb{Q}}$} isomorphes (i.e. $\exists$ isomorphisme de foncteurs $\alpha : f^{*} \simeq g^{*}$, tel que $\forall\, \overline{\eta}' \in \operatorname{Ob} \Pi_{K'}$, \uncertain{le} \struck{\ill{}} \add{\uncertain{carré}} \struck{correspondant}

\begin{tikzcd}
p f^{*}(\overline{\eta}') \arrow[r, "p(\alpha)", "\sim"'] \arrow[d, "\wr"'] & p g^{*}(\overline{\eta}') \arrow[d, "\wr"] \\
p'(\overline{\eta}') \arrow[r, "\sim"'] & p'(\overline{\eta}')
\end{tikzcd}

est commutatif) alors $f = g$.

\begin{tikzcd}
\Pi_{K'} \arrow[dr, "p'"'] & \overset{f^{*}}{\underset{g^{*}}{\rightrightarrows}} & \Pi_{K} \arrow[dl, "p"] \\
& \Pi_{\mathbb{Q}} &
\end{tikzcd}

L'hypothèse sur $f, g$ signifie aussi, en termes d'une clôture alg. choisie $\overline{K'}$ de $K'$ \struck{qui ne peut} \add{donnant via $f, g$ deux clôtures alg. de $K$, $\overline{K}$ et $\widetilde{K}$}, \struck{que les deux hom. d'exten\ill{}} que l'on peut trouver un isomorphisme entre celles-ci (\uncertain{permettant d'identifier} $\mathcal{E}_{\overline{K}/K}$ et $\mathcal{E}_{\widetilde{K}/K}$), \marginal{\uncertain{i.e.} induisant \uncertain{« l'identité »} sur les clôtures alg. $\overline{K}$ et $\widetilde{K}$ de \ill{}, identifiées à celles dans $\overline{K'}$} \struck{\ill{}} de telle \supplied{sorte} que les deux hom.
\[
f^{*}, g^{*} : \mathcal{E}_{\overline{K'}/K'} \rightrightarrows \mathcal{E}_{\overline{K}/K}
\]
soient égaux, \struck{Il \uncertain{s'agit donc des} \ill{}} \uncertain{Ceci s'exprime} \add{\uncertain{plus}} \ill{} en termes d'une clôture alg. $\overline{\mathbb{Q}}$ fixée de $\mathbb{Q}$, en disant que les deux hom. $f^{*}, g^{*} : \pi_{K'} \rightrightarrows \pi_{K}$ de groupes profinis extérieurs (avec opérateurs $\Gamma_{\overline{\mathbb{Q}}/\mathbb{Q}}$) sont égaux.

\marginal{On se ramène au cas $K = K'$ — c'est immédiat \uncertain{en} \struck{\uncertain{faisant le changement de base}} \ill{} \uncertain{catégories} finies \uncertain{sur} $\overline{\mathbb{Q}}$ \ill{} de $\mathbb{Q}$.}
\note{note marginale oblique en bas de la marge gauche ; en partie biffée.}

\page{22}
\note{pagination de l'auteur : 11.}

Écrivons comme d'habitude $K = \varinjlim A_i$, donc \struck{$K\otimes$} $\eta = \operatorname{Spec}(K) = \varprojlim U_i$, on a (en termes d'un pt géom. quelconque $\overline{\eta}$ de $\operatorname{Spec} K$ i.e. en termes d'un $\overline{K}$)
\[
\pi_K = \varprojlim_i \pi_1(\overline{U}_i, \overline{\eta}), \qquad \text{où } \overline{U}_i = U_i \otimes_{k} \overline{\mathbb{Q}}
\]
\note{l'indice du produit tensoriel est un $K$ ou un $k$ ; on lit $k$, le corps de base, ici $\mathbb{Q}$.}
et il suffit de voir que $\forall i$, $\underbrace{f \mid A_i}_{f_i} = \underbrace{g \mid A_i}_{g_i}$, moyennant le fait que
\[
\pi_1(f_i^{*}) = \pi_1(g_i^{*}) : \pi_{K'} \longrightarrow \pi_1(\overline{U}_i)
\]
(comme hom. de groupes extérieurs). Or pour $i$ fixé, on a $K' = \varinjlim A'_j$, où les $A'_j$ contiennent $f_i(A_i)$ et $g_i(A_i)$, donc
\[
\pi_{K'} = \varprojlim_j \pi_1(\overline{V}_j, \overline{\eta}'), \qquad \text{avec } \overline{V}_j = \operatorname{Spec}(A_j) \otimes_{k} \overline{\mathbb{Q}} .
\]
Notons (prenant les $V_j$ \struck{\uncertain{normaux}} \add{réguliers}) que les hom. de transition dans le syst. projectif des $\pi_1(\overline{V}_j, \overline{\eta}')$ sont surjectifs, donc $\pi_{K'} \to \pi_1(\overline{V}_j, \overline{\eta}')$ est surjectif, ce qui implique que l'égalité de $f^{*}$ et $g^{*} : \pi_{K'} \rightrightarrows \pi_1(\overline{U}_i)$ (comme hom. extérieurs) implique celle de $\pi_1(\overline{V}_j) \to \pi_1(\overline{U}_i)$. Donc l'égalité $f_i = g_i$ (d'où $f = g$) sera conséquence du résultat plus général,

\page{23}
\note{la page précédente s'achève sur « sera conséquence du résultat plus général, » ; ce résultat n'est pas énoncé en tête de cette page-ci, qui s'ouvre sur un corollaire.}

\struck{Corollaire} « purement géométrique » :

\textbf{Corollaire 1.} Soient $X, Y$ des schémas \add{de t.f. réduits \struck{\ill{}}} (sur un corps alg. clos $k$), et $f, g : X \to Y$ deux morphismes, on suppose que $\pi_1(f), \pi_1(g) : \pi_1(X) \rightrightarrows \pi_1(Y)$ sont égaux (en tant qu'hom. extérieurs). Alors

a) Si $Y$ \struck{est une extension d'une} \add{se plonge dans un $i : Y \to G$ (un groupe algébrique commutatif $G$)} \struck{\ill{} extension d'une} extension d'une VA par un tore, il existe un $a \in Y$ \uncertain{(unique)} tel que $g(x) = f(x) + a$ $\forall x \in X(k)$, i.e. $(i \circ g) = \tau_a \circ (i \circ f)$ ($\tau_a$ la translation par $a$).

\marginal{En fait, dans a) il suffit de supposer que $H_1(f, \mathbb{Z}_\ell) = H_1(g, \mathbb{Z}_\ell)$ pour un $\ell$ premier \uncertain{à la car.}}

b) Si $Y$ est une variété \uncertain{élémentaire} d'\uncertain{Artin} \ill{} \struck{\ill{}} fibrations successives des courbes \uncertain{anabéliennes}, \add{(et $X$ \struck{\uncertain{normal}} \add{réduit}) et si $f$ ou $g$ est dominant)} alors $f = g$.

\marginal{\uncertain{Un cas où} \ill{} \uncertain{premier} \ill{} \uncertain{structure des $\pi_1$} \ill{} \uncertain{l'unicité} \ill{} \uncertain{irréductibilité} \ill{}}
\note{note marginale oblique dans la marge gauche, en grande partie illisible, rattachée par un trait à l'énoncé b).}

\emph{Démonstration.} a) L'unicité de $a$ \struck{\ill{}} est claire — et \add{l'existence} \struck{est} \uncertain{revient}, (en choisissant \add{remplaçant $Y$ par $G$ et,} une « origine » $a \in X(k)$, et en se ramenant au cas $f(a) = g(a) = 0$, à prouver qu'alors $f = g$. D'ailleurs l'hyp. $\pi_1(f) = \pi_1(g)$ peut être affaiblie en l'hyp. que \struck{$f_1$ \ill{}} \add{\uncertain{$\ell$ premier à la car.}} $H_1(f, \mathbb{Z}_\ell) = H_1(g, \mathbb{Z}_\ell)$, — i.e. il suffit

\page{24}
\note{pagination de l'auteur : 12.}

d'examiner les actions de $\pi_1(f), \pi_1(g)$ sur les groupes abélianisés des $\pi_1$, et même sur leurs composantes $\ell$-adiques. \uncertain{Prenant alors} la jacobienne généralisée du type « ext. d'une V.A. par un tore » \struck{de} de $X$, on sait que

1°) les \struck{morphismes} $f : X \to G$ tels que $f(a) = 0$ se factorisent de façon unique par
\[
X \xrightarrow{\ \mathrm{can}\ } J \xrightarrow{\ \varphi\ } G,
\]
avec $\varphi$ un hom. de groupes algébriques

2°) \struck{$\varphi$ est} Un tel homomorphisme \uncertain{est} connu quand on connaît son action sur les $H_1(\ , \mathbb{Z}_\ell)$ (ce qui revient à le connaître sur les points d'ordre $\ell^\nu$, quel que soit $\nu$ — or ceux-ci sont denses …)

3°) $H_1(X, \mathbb{Z}_\ell) \simeq H_1(J, \mathbb{Z}_\ell)$

De ceci, \struck{\ill{}} on conclut (par 3°), que $H_1(f) = H_1(g)$ implique (vu $f = \varphi \circ \mathrm{can}$, $g = \psi \circ \mathrm{can}$) $H_1(\varphi) = H_1(\psi)$, donc par 2°) que $\varphi = \psi$, donc $f = g$, cqfd.

\struck{b)} Notons que l'on ne peut conclure de l'hyp. \add{dans a)} que $f = g$, puisque pour $X = Y = G$, les translations $\tau_a$ induisent l'identité sur $\pi_1(G)$ (du moins \struck{\ill{}} \uncertain{en car.} $0$ ou si $G$ est une VA), sans être l'identité.

b) On va pourtant prouver l'égalité,

\page{25}
sous l'hypothèse anabélienne faite, et celle de dominance. (Cette dernière n'est pas superflue, comme on voit en prenant \struck{$Y = G \times Z$, où $G$ \ill{} $Y$ un groupe elliptique et $Z$ un sous-groupe fini $\neq 1$, et pour $f, g$} \add{deux applications constantes de $X$ dans $Y$ !})

Tout d'abord, on vérifie qu'une variété élémentaire d'Artin se plonge \add{par un $i : Y \to G$} \add{\uncertain{bien sûr}} dans une ext. d'une V.A. par un tore — donc on aura $(i \circ g) = \tau_u (i \circ f)$. \struck{Donc} \add{D'autre part}, \struck{Quitte à remplacer} $f, g$ étant dominants, il existe une partie ouverte \add{non vide} $V \subset Y$ contenue dans $\operatorname{Im} f$, et \ill{} $V + u \subset Y$. Soit $U \subset X$ un ouvert \add{\uncertain{dense}} tel que $f(U) \subset V$ \struck{$g(U) \subset V$}, \struck{donc \ill{}} quitte à remplacer $X$ par $U$, on \ill{} se ramène au cas où $f$ \struck{et $g$} envoie $X$ dans un ouvert $V$ de $Y$ tel que $V + u \subset Y$, où $g \circ i = \tau_u (i \circ f)$. On \struck{\ill{}} \add{voit} que $\pi_1(f)$ est un composé (où $f' : X \to U$ est induit par $f$ et $j : U \to Y$ est l'inclusion)
\[
\pi_1(X) \xrightarrow{\ \pi_1(f')\ } \pi_1(U) \xrightarrow{\ \pi_1(j)\ } \pi_1(Y)
\]
tandis que $\pi_1(g)$ est le composé
\[
\pi_1(X) \xrightarrow{\ \pi_1(f')\ } \pi_1(U) \xrightarrow{\ \pi_1(j')\ } \pi_1(Y)
\]
où $j' : U \to Y$ est induit par $\tau_u$.
\note{dans ces dernières lignes l'ouvert de $Y$, d'abord noté $V$, est noté $U$ comme l'ouvert de $X$ ; on transcrit tel quel.}

\page{26}
\note{pagination de l'auteur : 13.}

L'hyp. que $\pi_1(f)$ \struck{et} $= \pi_1(g)$ \emph{signifie} donc qu'il existe $\alpha \in \pi_1(Y)$, tel que
\[
\pi_1(j')(\gamma) = \operatorname{int}(\alpha)\, \pi_1(j)(\gamma)
\]
pour tt $\gamma \in \operatorname{Im}\bigl(\pi_1(X) \xrightarrow{\pi_1(f)} \pi_1(U)\bigr)$. Or cette image est un sous-groupe \emph{ouvert} de $\pi_1(U)$ (\uncertain{fait} vrai pour un morphisme dominant !). \struck{D'ailleurs} Donc on est ramené à ceci : Soit $U$ ouvert \add{$\neq \emptyset$} de $Y$, $u \in G$, tels que $\tau_u U \subset Y$ [et tels que (désignant par $j, j'$ les \struck{applications} \add{morphismes} $y \mapsto y$ et $y \mapsto y + u$ de $U$ dans $Y$) $\pi_1(j)$ et $\pi_1(j')$ coïncident extérieurement sur un sous-groupe ouvert de $\pi_1(U)$] — \emph{alors} $j = j'$ i.e. $u = 0$.

\marginal{NB ceci n'est pas plus fort que le corollaire, comme on voit en prenant pour $X$ le \uncertain{rev. ét. fini} de $U$ correspondant à un sous-groupe ouvert \uncertain{donné} de $\pi_1(U)$ \ill{}}

Finalement, je n'arrive pas à le prouver par voie \emph{géométrique} — pour ceci je me rabats sur une méthode \emph{arithmétique},
\struck{ce qui fournit d'ailleurs \struck{affaibli} généralisé le th. géométrique ; on notera que \ill{} remplacer $X$ par \ill{} ouvert non vide convenable \ill{} \ill{} par les \ill{} \ill{} \ill{} que $X$ est une variété}
\marginal{Si : cf. plus bas !}
\note{le passage biffé (cadre et traits obliques) se poursuit en tête de la page suivante ; la note marginale « Si : cf. plus bas ! » est écrite contre lui.}

\page{27}
\struck{élémentaire d'Artin \add{anabélienne} — donc sans $\pi_1$ \struck{\ill{}} — \uncertain{à} centre trivial. Ceci posé}
\note{fin du passage biffé commencé à la page précédente.}

\textbf{Corollaire 2.} La conclusion \add{$f = g$} du corollaire précédent \struck{b)}, \struck{est} reste valable si \struck{au lieu de supposer $f$ ou $g$ dominant} on suppose \add{que $k$ est de car. $0$,} $X$ \emph{normal}, et \struck{satisfaite} dans l'une des hyp. suivantes :

c) l'image de $\pi_1(f)$ est un sous-groupe ouvert de $\pi_1(Y)$, $Y$ est une variété élémentaire d'Artin anabélienne

d) l'image \add{de $\pi_1(X)$} par $\pi_1(f)$ \struck{de tout sous-groupe ouvert de $\pi_1(X)$} a un centralisateur dans $\pi_1(Y)$ réduit à $(1)$, \struck{et $X$} \add{et $Y$} se plonge dans un groupe alg. extension d'une VA par un tore.

\struck{e) Le centre de $\pi_1(X)$, et les centralisateurs de l'image dans $\pi_1(Y)$, sont réduits à $1$.}
\marginal{\struck{\ill{}}}
\note{l'alinéa e) est annulé par une série de boucles ; une note marginale encadrée, biffée, s'y rattache.}

Comme les centralisateurs d'un sous-groupe ouvert de $\pi_1(Y)$ ($\pi_1(Y)$ étant ext. successive de groupes profinis \emph{libres} anabéliens) est réduit à $(1)$, comme on a vu${}^{*}$ plus haut, le cas c) est un cas particulier de d), \struck{des variétés $X$ \ill{}} \struck{D'autre part}, \struck{— mais d) se ramène à e), en remplaçant $X$ par un ouvert d'Artin} (dans le cas d), quitte à remplacer $X$ par un ouvert d'Artin (la normalité assurant que $\pi_1(U) \to \pi_1(X)$ surj.))

\marginal{${}^{*}$ il faut quand même traiter à part le cas des $\pi_1$ d'une courbe complète de genre $g \geq 2$, et vérifier que son $\pi_1$ est à centre trivial. \ill{} \uncertain{d'Artin} \ill{} \uncertain{de borner} \ill{} \uncertain{dans le cas des fibrations successives} \ill{}}

\page{28}
\struck{OPS de plus que $\pi_1(X)$ a un centre trivial. En car. $p > 0$, \uncertain{il suffira} de travailler avec les groupes fondamentaux premiers à $p$ — \ill{} sur les localisés en $\ell \neq p$ des groupes fondamentaux.}
\note{ces lignes, en tête de page, sont biffées par des traits obliques et un trait horizontal ; la pagination de l'auteur, 14, est écrite au-dessous d'elles.}

La situation $X, Y, f, g$ provient, par extension du corps de base, d'une situation analogue sur un corps $K$ \add{extension} (de type fini \struck{de} de $\mathbb{Q}$). \struck{Choisissons une} \add{Soit $\overline{K}$ la} clôture algébrique \struck{$\overline{\mathbb{Q}}$} de $K$ dans $k$, \struck{\ill{}} et $\overline{X}, \overline{Y}, \overline{f}, \overline{g}$ \struck{sur} déduits de \add{$X, Y,$} $f, g$ par extension du corps de base de $K$ à $\overline{K}$. On a donc $\pi_1(\overline{f}), \pi_1(\overline{g})$ satisfaisant la condition d) \uncertain{et} de plus $\pi_1(\overline{X})$ a centre trivial.

Mais nos hypothèses impliquent que les extensions $\mathcal{E}_{X/K}$ \add{$= \pi_1(X)$}, $\mathcal{E}_{Y/K}$ \add{$= \pi_1(Y)$} \add{de} $\mathcal{E}_{\overline{K}/K}$ par $\pi_1(\overline{X})$, resp. $\pi_1(\overline{Y})$, ainsi que les hom. $\pi_1(f), \pi_1(g)$ induits, sont reconstruits à partir \add{de $\pi_1(\overline{X}), \pi_1(\overline{Y})$, de} $\pi_1(\overline{f}), \pi_1(\overline{g})$ et de l'action extérieure de $\mathcal{E}_{\overline{K}/K}$ sur ces groupes. On va montrer maintenant le

\page{29}
\textbf{Corollaire 3.} Soient \struck{$f, g$} $X, Y$ deux schémas de type fini sur un corps $K$ ext. de type fini de $\mathbb{Q}$, \add{on suppose} \struck{tels} que $Y$ se plonge dans un \struck{VA} extension d'une V.A. par un tore, $X$ réduit, \add{$X, Y$ géom. connexes.} \struck{d'où} Soit $\overline{K}$ une clôture algébrique de $K$, d'où des extensions « extérieures » $\mathcal{E}_{X/K}$, $\mathcal{E}_{Y/K}$ \add{de $\mathcal{E}_{\overline{K}/K}$} $= \operatorname{Gal}(\overline{K}/K)$ par $\pi_1(\overline{X}), \pi_1(\overline{Y})$, et pour tt morphisme $f : X \to Y$, un morphisme $\pi_1(f)$ de $\mathcal{E}_{X/K}$ \struck{dans} dans $\mathcal{E}_{Y/K}$. Soient $f, g : X \rightrightarrows Y$ tels que $\pi_1(f) = \pi_1(g)$ — i.e. $\pi_1(f)$ et $\pi_1(g)$ soient extérieurement conjugués, par un élément de $\pi_1(Y)$ — alors $f = g$.

En fait, il suffit \struck{\ill{}} que les hom. \add{donc d'extensions tout court !)} d'\emph{extérieures} (extérieures, par $\pi_1(\overline{X})_{\mathrm{ab}}$ et par $\pi_1(\overline{Y})_{\mathrm{ab}}$ soient égaux, pour pouvoir conclure $f = g$. (C'est à dire, une fois que, par des hypothèses \add{géométriques} \emph{anabéliennes}, des hypothèses purement géométriques sur les actions de $\pi_1(\overline{f}), \pi_1(\overline{g})$ sont par de
\marginal{\uncertain{i.e.} \ill{} $H_1(\overline{X}, \mathbb{Z}_\ell)$ \ill{} $H_1(f, \mathbb{Z}_\ell)$}

\page{30}
\note{pagination de l'auteur : 15.}

ramenées à des hypothèses diverses \add{de nature} arithmético-géométrique, on peut laisser tomber les aspects anabéliens, pour jouer \uncertain{à fond} sur les aspects abéliens — mais dans un contexte arithmétique. Au lieu d'avoir \add{et d'utiliser} \struck{\ill{}}, comme dans le cor. 1, a), l'hypothèse que $H_1(\overline{f}) = H_1(\overline{g}) ; H_1(\overline{X}) \to H_1(\overline{Y})$, on a une hypothèse plus fine, sur l'égalité de deux homomorphismes d'\emph{extensions} de $\mathcal{E}_{\overline{K}/K}$ par $H_1(\overline{X}, \widehat{\mathbb{Z}}_\ell)$ resp. $H_1(\overline{Y}, \widehat{\mathbb{Z}}_\ell)$ $\forall \ell \in \mathbb{P}$.)
\note{les coefficients sont écrits $\mathbb{Z}$ coiffé d'un chapeau, avec un indice qui paraît être $\ell$.}

Démonstration du corollaire 3.

Il suffit de voir que $f(x) = g(x)$ pour tt $x \in X(K)$ (car, quitte à faire une extension finie \struck{\ill{}} du corps de base, cela \struck{reste} vrai pour $x \in X(\overline{K})$). \struck{puis \ill{} extension des corps}

Or un point définit une classe de \struck{d'autre part} \add{conjugaison de} scindages de $\mathcal{E}_{X/K}$ sur $\mathcal{E}_{\overline{K}/K}$, et a fortiori de l'extension « abélianisée » associée — l'hypothèse implique que $f(x)$ et $g(x)$ définissent la même classe de

\page{31}
conjugaison de scindages. Donc il suffit maintenant de prouver le

\textbf{Théorème 3.} Soit $X$ un schéma de type fini sur un corps $K$, extension de type fini de $\mathbb{Q}$, on suppose que $X$ est \struck{connexe lisse} \add{géom. connexe, et} se plonge dans une extension d'une VA par un tore (\uncertain{p. ex.} $X$ est une variété élémentaire d'Artin, à fibres \uncertain{types} \add{\uncertain{génériques}} distinctes de la droite projective et de la droite affine). Considérons \add{une} \struck{l'application} clôture alg. $\overline{K}/K$, et l'extension extérieure \add{correspondante} $\mathcal{E}_{X/K}$ \add{de} $\mathcal{E}_{\overline{K}/K} = \operatorname{Gal}(\overline{K}/K)$ par $\pi_1(\overline{X})$ ($\overline{X} = X \otimes_K \overline{K}$), et l'extension $\widetilde{\mathcal{E}}_{X/K}$ de $\mathcal{E}_{\overline{K}/K}$ par $\pi_1(\overline{X})_{\mathrm{ab}}$ \uncertain{déduite}. Considérons les applications canoniques
\begin{quote}
(*) $X(K) \longrightarrow$ classes de \add{$\pi(\overline{X})$-}conjugaison de scindages de $\mathcal{E}_{X/K}$ sur $\mathcal{E}_{\overline{K}/K}$

(**) $X(K) \longrightarrow$ classes de conjugaison de scindages de l'ext. $\widetilde{\mathcal{E}}_{X/K}$ de $\mathcal{E}_{\overline{K}/K}$ par $H_1(\overline{X})$.
\end{quote}
\note{entre les deux lignes, une flèche pointillée verticale descend de la première vers la seconde.}

\marginal{il suffirait $K$ de type fini sur son corps premier, pas néc. de car. $0$ — mais sous cette forme \uncertain{est} pas très utilisable en car. $p > 0$, cf. \ill{}}

\page{32}
\note{pagination de l'auteur : 16.}

Ces applications sont \emph{injectives}.

\emph{Démonstration.} Il suffit de le prouver pour la seconde application, et on est ramené au cas où $X$ est \add{lui-même} un groupe alg. \add{$G$}, ext. d'une VA par un tore. Alors l'application est un hom. de groupes
\[
(16) \qquad \underset{H^0(K, G)}{G(K)} \longrightarrow H^1(K, H_1(\overline{G}))
\]
obtenu ainsi. On considère pour tt entier $n$ la suite exacte kummérienne
\[
0 \to {}_nG \to G \xrightarrow{\ n\,\mathrm{id}_G\ } G \to 1
\]
d'où suite exacte de cohomologie, qui donne
\[
0 \to H^0(K, G)_n \to H^1(K, {}_nG) \to {}_nH^1(K, G) \to 0
\]
et passant à la limite, on trouve
\[
0 \to \varprojlim_n H^0(K, G)_n \to \varprojlim_n H^1(K, {}_nG) \to \varprojlim_n {}_nH^1(K, G)
\]
Le composé de (16) avec l'hom. \struck{canonique}
\[
H^1(K, H_1(\overline{G})) \to \varprojlim_n H^1(K, {}_nG)
\]
compte tenu de
\[
H_1(\overline{G}) = \varprojlim_n {}_nG
\]
montre que l'homomorphisme induit par
\[
H^0(K) \longrightarrow \varprojlim H^0(K, G)_n \hookrightarrow \varprojlim H^1(K, {}_nG)
\]

\page{33}
dont le noyau \add{est $\bigcap_n nH^0(K, G)$, donc} est formé des éléments de $G(K)$ \emph{infiniment divisibles} dans $G$. Or ici $K$ étant un \struck{groupe} corps \add{absolument} de type fini, le th. de Mordell-Weil nous assure que $G(K)$ est un \struck{groupe} \add{$\mathbb{Z}$-module de} type fini — donc $G(K) \to \varprojlim G(K)_n$ est injectif. Donc $G(K) \to \varprojlim H^1(K, {}_nG)$ est injectif, et a fortiori l'hom (16), cqfd.

\emph{Remarque.} Il ne suffisait pas dans cet argument de prendre, pour un $\ell$, l'extension de $\mathcal{E}_{\overline{K}/K}$ par $H_1(X, \mathbb{Z}_\ell)$ — il faut utiliser tous les $\ell$ — sinon, \ill{} $p$, alors la translation par un point d'ordre $p$ donne un contre-exemple. \struck{\ill{}} Le résultat revient à ceci : Pour tt $x \in G(K)$, on regarde le torseur sur $T_\infty(G) = \varprojlim {}_nG$, image \uncertain{sur} $K$, de groupe

\page{34}
\note{pagination de l'auteur : 17.}

inverse de $x$ dans le « rev. universel abélien » $\widetilde{G}$ de $G$, construit comme $\varprojlim$ des revêtements $G(n) \simeq G$ de $G$, donnés par $n\,\mathrm{id}$. L'énoncé dit que si ce torseur est trivial — i.e. si $\forall n$, $x$ \uncertain{se remonte} — alors $x = 0$. C'est vrai en car. $p > 0$ également — mais dans ce cas les ${}_nG$ ne sont pas étales. Si on suppose \ill{} \uncertain{qu'on remplace} les ${}_pG$ par leurs \uncertain{quotients} \ill{} pour leurs parties radicielles — on travaille dans la cat. des « quasi-variétés » où on \add{\uncertain{inverse}} les \struck{\ill{}} \add{\uncertain{morphismes}} finis radiciels surjectifs, \struck{identifie} \struck{\ill{} il} \ill{} \uncertain{avec la partie première} \ill{} $H_1(\overline{G})$, \uncertain{qui dit} l'existence \add{des relèvements} \ill{} radiciels $/K$ implique déjà $x = 0$. (\ill{} \add{Bien sûr $x$ est de $p$-torsion,} \ill{} $px = 0$ ….) Mais \uncertain{ramené au cas d'une VA}, \ill{} la validité sous cette forme du th. 3 en car. $p > 0$ \uncertain{me} fait une belle jambe — en car. $p > 0$ \ill{}, pour être à l'aise, il faudrait pouvoir le prouver pour les contextes des groupes fondamentaux « premiers à $p$ ». Or déjà dans le cas d'une courbe elliptique $Y$ \add{« contextes »} \ill{} \uncertain{(en se bornant aux \ill{} des pts)} \struck{\ill{}} \uncertain{l'injectivité de} $(**)$ est \emph{fausse}. Il

\marginal{O.K. \ill{} \uncertain{remplaçant} \ill{} $K$ \ill{} \uncertain{radiciels} \ill{}}
\note{deux notes marginales obliques, l'une biffée en tête de la marge gauche, l'autre commençant par « O.K. » vers le bas, sont pour l'essentiel illisibles.}

\page{35}
est cependant possible que l'injectivité de $(*)$ reste valable sans hyp. de caractéristique — à examiner de plus près…

Revenant aux conditions du cor. 1, b), en car. $p > 0$, et travaillant avec les groupes fondamentaux premiers à $p$, on trouve que le $u$ dans 1°) est un point \add{de $G(K)$} d'ordre une puissance de $p$. Si $G$ est un tore \ill{} pour celui que les « fibres » dans la variété élémentaire $Y$ soient de genre $0$) on en conclut encore que $u = 0$, i.e. $f = g$. Dans le cas général, la démonstration arithmétique prouve \add{(au moins)} ceci : si $x \in X$, \add{d'où} $f(x) = P$, $g(x) = Q$, \uncertain{on voit} que $Q - P$ est un élément de $p$-torsion de la jacobienne génér. $G_Y$ de $Y$, \add{de degré premier à $p$,} \ill{} pour tt rev. \add{étale} galoisien fini $Y'$ de $X$, \uncertain{et} il existe pour un $P'$ au-dessus de $P$, \uncertain{il existe} $Q'$ au-dessus de $Q$ tel que $Q' - P'$ soit un élément de $p$-torsion de $G_{Y'}$, \struck{\ill{}} \add{\ill{} divisible par les} \struck{comme \ill{} relèvements les isogénies $n_G$ ($n$ premier à $G$), il existe un $Q_n \in G$ tel que $n \cdot Q_n = Q$, et que $Q_n$} Il se pourrait qu'on puisse en déduire $P = Q$.

\marginal{Certainement le th. 2 reste vrai sur un corps de base de car. $p > 0$, en travaillant avec les groupes fondamentaux premiers à $p$ — car \uncertain{dans} la réduction à un énoncé géométrique \ill{} la suite \uncertain{génér.} \ill{} \uncertain{peut supposer} que $Y$ est un \emph{produit} de courbes de genre zéro. Dans ce cas d'ailleurs, on est ramené au cas où $Y$ est une telle courbe ….}

\page{36}
\note{pagination de l'auteur : 18.}

\emph{Complément.} Retour sur une démonstration \emph{géométrique} du th. 2, cor. 1 b). OPS que $G$ est la jacobienne généralisée de $Y$, et il suffit de montrer le

\emph{Lemme.} Soit $Y$ une variété élémentaire d'Artin anabélienne (sur $k$ alg. clos), $Y \hookrightarrow J^1_Y$ son plongement dans sa jacobienne généralisée, $u \in J^0_Y(k)$ et $U$ un ouvert non vide de $Y$, tels que $U + u \subset Y$. Alors $u = 1$, \add{en d'autres termes : l'application $x \mapsto x + u$ de $U$ dans $Y$ est l'identité.}
\note{l'énoncé du lemme est marqué d'un double trait dans la marge.}

Par dévissage, on est ramené au cas où $Y$ est une courbe. Supposons-la d'abord complète, de sorte que $U + u \subset Y$ implique $Y + u \subset Y$ — alors le résultat est bien connu : \uncertain{c'est} \ill{} p. ex. de la formule des pts fixes, qui implique que \add{si \ill{} n'est l'identité,} $\tau_u \mid Y$, a $2 - 2g$ pts fixes, ce qui (comme $g \geq 2$ en vertu de l'hypothèse anabélienne) est absurde \add{au moins si $g \neq 1$,}.

\marginal{Si $g \geq 2$, parfait : remplaçant $Y$ par $\hat{Y}$. OPS $Y$ complète. Si $g = 1$ ou $0$, cf. arguments cohomologiques plus bas}

Cela \struck{signifie} \add{implique} donc, dans le cas général, que l'image de $u$ dans la partie abélienne de $J^0_Y(k)$ est nulle. \struck{[[Or d'après la structure connue de $J^0_Y$ comme extension de $J^0_{\hat{Y}}$ par un tore $\mathbb{G}_m^I/\mathbb{G}_m$ (où $I = \hat{Y} \setminus Y$), $u$ s'identifie à un élément de $(\mathbb{G}_m^I/\mathbb{G}_m)(k)$, donc est défini par $(\lambda_i)_{i \in I}$, les $\lambda_i \in k^{*}$, qui servent à construire un fibré inversible}
\marginal{on \uncertain{peut} voir ici :}
\note{le passage ouvert par « [[ » est barré d'un trait oblique ; le biffage se poursuit page suivante.}

\page{37}
\struck{sur la courbe « pincée » $Y' = \hat{Y}/I$ (déduite de $\hat{Y}$ en contractant $I$ en un pt) en identifiant les fibres \add{réduites en les $y_i \in I$} du faisceau trivial sur $\hat{Y}$, \struck{\ill{}} à $k$ par les $\lambda_i$ — entre elles. Alors $u + x$ (pour $x \in Y$) est représenté par le faisceau $\mathcal{O}_{\hat{Y}}[x]$ (associé au diviseur $x$), \add{de $J^1_{\hat{Y}}$} \uncertain{muni} de ces mêmes identifications.]}
\note{fin du passage biffé (trait oblique) ouvert page précédente ; il est encadré d'un trait vertical et fermé par « ] ».}

Pour que $x + u$ \struck{\ill{}} soit de la forme $y$ ($y \in Y$) il faut donc que $u + x$ et $y$ aient même image dans $J^1_{\hat{Y}}$ \add{et comme l'image de $u$ dans $J^0_{\hat{Y}}$ est nulle}, ce que \struck{(pour un genre $g \geq 1$)} (pour un $g \geq 1$) implique déjà $x = y$, donc $u = 0$. Si $g = 0$, \add{\ill{}} donc $\nu = \operatorname{card} I \geq 3$ en vertu de l'hyp. anabélienne, on \struck{remplace} prend une partie $J$ de $I$ à deux éléments, et on remplace $\hat{Y}$ par \uncertain{$\hat{Y}_J$}] pour conclure encore à $x = y$ donc $u = 0$.

\marginal{Conclusion : on saurait \ill{} que l'image de $u$ dans $J^0_{\hat{Y}}$ est nulle !}

[En fait, on a prouvé que si $u \in \operatorname{Ker}(J^0_Y \to J^0_{\hat{Y}})$, alors $(Y + u) \cap Y = \emptyset$ — et par ailleurs, s'il existe un ouvert non vide $U$ de $Y$ tel que $U + u \subset Y$, alors $u$ est dans le Ker précédent \ill{} $\neq 1$). Mais ces arguments ne marchent que pour $g \neq 1$.] Il vaut mieux, dans le cas général, présenter les choses sous forme cohomologique.

\marginal{$g \geq 2$ ?}

Considérons les deux \struck{applications} \add{morphismes} $U \rightrightarrows Y$, inclusion

\page{38}
\note{pagination de l'auteur : 19.}

et $j : U \to Y$ induit par $\tau_u$, je dis que $H_1(i) = H_1(j)$, ou ce qui revient au même, puisque $Y \xrightarrow{\alpha} J^1_Y$ induit un iso. $H_1(\alpha) : H_1(Y) \to H_1(J^1_Y)$, que $H_1(\alpha) H_1(i) = H_1(\alpha) H_1(j)$, ou encore $H_1(\alpha i) = H_1(\alpha j)$. Or on a $\alpha j = \tau_u (\alpha i)$ donc \struck{il suffit de}
\[
H_1(\alpha j) = H_1(\tau_u)\, H_1(\alpha i)
\]
et on sait que $H_1(\tau_u) = \mathrm{id}$, OK. Cela étant, soit $I = \hat{Y} \setminus Y$, $J = \hat{Y} \setminus U$, de telle façon que $J \supset I$, et soit $f : \hat{Y} \simeq \hat{Y}$ déduit de $j : U \hookrightarrow Y$ en passant aux complétés, donc \struck{on a} $f(U) = \hat{Y} - f(J) \subset Y = \hat{Y} - I$, \uncertain{implique} que $f(J) \supset I$ i.e. $J \supset f^{-1}(I)$. On voit que le \uncertain{noyau} de $H_1(U) \to H_1(Y)$, \ill{} \ill{} en $J \setminus I$, \ill{} \add{\ill{}} $L_j$ avec $j \in J \setminus I$, et pour le noyau analogue de $H_1(U) \to H_1(Y)$ induit par $j$ est le \ill{} des $L_j$ avec $j \in J \setminus f(I)$, d'ailleurs les $L_j$ ($j \in J$) dans $H_1(U)$ sont $\simeq \widehat{\mathbb{Z}}$ (resp. $\widehat{\mathbb{Z}}'$ en car. $p > 0$ …) et le sous-$\widehat{\mathbb{Z}}$-module engendré est
\[
\widehat{\mathbb{Z}}^J / \widehat{\mathbb{Z}} \ (\text{diagonale})
\]
, donc \struck{si $\operatorname{card} J \geq 3$,} les sous-modules engendrés par deux parties \add{$I, I'$} de $J$ distinctes de $J$, ne peuvent être égaux que si $I = I'$, ou $I$ et $I'$ de cardinal $\mu - 1$ (où $\mu = \operatorname{card} J$). Mais ici, quitte à rapetisser $U$, on peut supposer

\page{39}
que $\operatorname{card} I < \operatorname{card} J - 1$ i.e. $\operatorname{card}(J \setminus I) = \operatorname{card}(Y \setminus U) > 1$. \struck{On est donc dans l'un des cas suivants}

\struck{a) $f^{-1}(I) = I$ i.e. $I$ stable par $f$.}

\struck{b)} Donc on voit que l'on doit avoir $f^{-1}(I) = I$ i.e. $I$ stable par $f$, ce qui montre que $\tau_u$ induit en fait un \emph{automorphisme} de $Y$. \add{Remplaçant} \struck{\ill{}}

\struck{$Y$ est une courbe elliptique} dans ces arguments $Y = \hat{Y} \setminus I$ par $Y' = Y \cup i = \hat{Y} \setminus I'$ ($I' = I - \{i\}$), on voit de même que $f$ induit un automorphisme de $Y'$ (cet argument \add{\uncertain{montre}} que $Y'$ s'envoie dans \uncertain{sa} jacobienne généralisée \uncertain{ce qui} \ill{} anabélienne sur $Y$) \struck{\ill{} le cas où $Y$ est de type $(0, \nu)$} \add{\uncertain{exclu par l'hyp. anabélienne}} \struck{\ill{}} donc $f$ induit l'identité sur $I$. Si le genre est $0$, on en conclut (puisque $\nu \geq 3$) que $f$ est l'identité ; \struck{Reste le} Dans le cas du genre $1$, on en conclut maintenant que l'image de $u$ dans $J^0_{\hat{Y}}$ est égale à $1$, et on termine comme précédemment.

Ouf !

\page{40}
\note{pagination de l'auteur : 20.}

\textbf{2. La question de pleine fidélité.}

Soient $K, K'$ deux extensions de type fini de $\mathbb{Q}$ — est-il vrai que tout $\Pi_{\mathbb{Q}}$-homomorphisme \uncertain{$\Pi_{K'} \to \Pi_{K}$} provient d'un homomorphisme de corps $K' \to K$ ?
\note{les indices de $\Pi_{K'} \to \Pi_K$ se lisent mal ; plus bas la question porte sur les hom. de $\pi_{K'}$ dans $\pi_K$ induits par un hom. $K \hookrightarrow K'$.}
On est ramené aussitôt au cas où — une clôture algébrique $\overline{\mathbb{Q}}$ de $\mathbb{Q}$ étant choisie, d'où un $\Gamma_{\overline{\mathbb{Q}}/\mathbb{Q}}$ — $K$ et $K'$ ont des sous-corps $k, k'$ (clôture alg. de $\mathbb{Q}$ dans $K$ resp. $K'$) isomorphes, avec des \struck{iso} \add{plongements} $k, k' \to \overline{\mathbb{Q}}$ de même image, \struck{donc} \uncertain{de sorte} que $\mathcal{E}_K$ et $\mathcal{E}_{K'}$ peuvent être considérées comme des extensions d'un même groupe $\Gamma = \Gamma_{\overline{\mathbb{Q}}/k}$ par $\pi_K$ resp. $\pi_{K'}$. La question est alors si \emph{tout} hom. de $\pi_{K'}$ dans $\pi_K$, qui commute à l'action de $\Gamma$, est induit par un hom. $K \hookrightarrow K'$. Pour construire un tel hom., il faudrait donc avoir une idée comment reconstruire $K, K'$ à partir des extensions $\mathcal{E}_K, \mathcal{E}_{K'}$, — ou encore à partir
\note{la phrase se poursuit au-delà de ce lot.}

\end{document}
